\chapter{Introduction}
\label{ch:introduction}

\section{What is BioDYM?}

BioDYM is a Material Flow Analysis (MFA) framework built on top of the Open Dynamic Material Systems Model (ODYM). It provides a complete workflow for modeling material flows, stocks, and transformations in bio-based and industrial systems.

BioDYM extends ODYM with:
\begin{itemize}
  \item \textbf{Dynamic Stock Modeling (DSM)} --- Cohort-based stock tracking with configurable lifetime distributions (Normal, Weibull, LogNormal, etc.)
  \item \textbf{First-Order Multi-Pool decay (FOMP)} --- Two-pool organic matter decomposition (labile and recalcitrant fractions)
  \item \textbf{Monte Carlo Simulation} --- Uncertainty quantification with multiple distribution types and sensitivity analysis
  \item \textbf{Scenario Analysis} --- Define and compare parameter modifications across multiple scenarios
  \item \textbf{Excel-based Configuration} --- All model inputs defined in a single Excel workbook
  \item \textbf{Interactive Visualization} --- Sankey diagrams, dynamics plots, stock analysis, and publication-quality export
\end{itemize}

\section{The Element Concept}

BioDYM tracks multiple \emph{elements} simultaneously across all flows and stocks. Elements represent mass-conserved quantities that are tracked in parallel:

\begin{table}[h]
\centering
\begin{tabular}{llp{7cm}}
\toprule
\textbf{Element} & \textbf{Abbreviation} & \textbf{Description} \\
\midrule
Material & material & Total mass (always the first element) \\
Water Content & WC & Mass of water in the flow \\
Dry Matter & DM & Mass of dry matter ($\text{material} = \text{WC} + \text{DM}$) \\
Carbon Content & CC & Mass of carbon in the dry matter \\
\bottomrule
\end{tabular}
\caption{Default element configuration for biomass systems.}
\label{tab:elements}
\end{table}

Elements can be configured for any application. For example, a metals system might use \texttt{[material, Fe, Cu, Al]}, while a packaging system might use \texttt{[material, PET, PE, PP]}.

\subsection{Element Hierarchy}

Elements can have parent--child relationships. In the default biomass configuration:
\begin{itemize}
  \item DM is defined as a \emph{percentage of material}
  \item CC is defined as a \emph{percentage of DM} (not of material)
\end{itemize}

This hierarchy is configured in the Excel workbook and automatically maintained by the calculation engine.

\section{System Architecture}

A BioDYM model consists of:
\begin{description}
  \item[Processes] Nodes in the system (e.g., Forestry, Processing, Use Phase)
  \item[Flows] Directed connections between processes, carrying mass of each element over time
  \item[Stocks] Material accumulated within a process
  \item[Transfer Coefficients (TCs)] Parameters controlling how material is distributed from a process to its outflows
  \item[Parameters] All model inputs: flow data, compositions, TC values, DSM/FOMP configurations
\end{description}

Process~0 is always the \emph{system boundary} (environment). Flows from Process~0 into the system are inputs; flows to Process~0 are outputs.

\section{Process Types}

Each process has a \emph{logic} that determines how it handles material:

\begin{table}[h]
\centering
\begin{tabular}{lp{9cm}}
\toprule
\textbf{Process Logic} & \textbf{Behavior} \\
\midrule
Input & System entry point (boundary process) \\
Splitter & Distributes material to outflows using TCs; preserves element composition \\
Transformer & Distributes material using element-specific TCs; changes composition \\
Pass-through & Passes all inflow directly to a single outflow \\
DSM & Dynamic stock with cohort-based lifetime tracking \\
FOMP & First-order multi-pool organic matter decay \\
\bottomrule
\end{tabular}
\caption{Available process logic types in BioDYM.}
\label{tab:process-types}
\end{table}
